\subsection*{The Convention: Allowed Argument Types}

In current implementation, the argument must be either:

\begin{itemize}[leftmargin=1.6em]
\item \emph{Spanned Data} type, or
\item Scalar \emph{Integer} type, or
\item Scalar \emph{Floating-Point} type.
\end{itemize}

Here, the scalar \emph{integer} type only includes \texttt{int}, whereas the scalar \emph{floating-point type} includes \texttt{float}, \texttt{double}, \texttt{half}, \texttt{bfp16}, and \texttt{half8}. However, since the target-platform differs in floating-point support, the \emph{floating-point type} parameter also varies for different target platforms.

In addition, \textbf{operations on floating-point are not supported} by Croqtile till now, considering the floating is not useful for tileflow programs. Instead, they are typically values for device computation, which the \emph{tileflow program}s pass such values directly to device kernels.

The below code showcases an example:

\begin{lstlisting}[language=croqtile,caption={The Convention: Allowed Argument Types}]
__device__ void bar(float *p, int m, int n) {}
__device__ void foo(float *, int, unsigned, double, float) {}
__co__ void foobar(f32 [M, 24] input, int N, float padding) {
  parallel p by 1 {
    shared f32 [14, 7] buffer;
    call bar(input.data, M, buffer.span(0));
    call foo(buffer.data, N, 3, 3.14, padding);
  }
}
\end{lstlisting}

In this example, it passes different data types as arguments from Croqtile function to the device functions, which matches the device parameters exactly. Note that in the device function, the **corresponding parameter type of the \emph{spanned data} argument is simply the pointer of its \_element type\_**, where the shape information is dropped. For example, the \texttt{foo}'s argument \texttt{p} is \texttt{float*}, which corresponds the \emph{spanned data} argument \texttt{input.data}. In Croqtile, you names the parameter as a \textbf{decayed} pointer of the \emph{spanned data}. If using other form like \texttt{float p[]}, it will trigger failure in target compilation stage.

For some types like \texttt{f16}, and \texttt{bf16}, there may not be native target support of such types, it is possible to utilize \texttt{croqtile::f16} and \texttt{croqtile::bf16} to handle such types.
